\documentclass[11pt,a4paper]{article}
\usepackage{geometry}
\geometry{margin=2.5cm}
\usepackage{enumitem}
\setlist{topsep=2pt,itemsep=2pt}
\usepackage{booktabs}
\usepackage{array}
\usepackage{amsmath,amssymb}
\usepackage{graphicx}
\usepackage{subcaption}

\title{Bayesian Workflow -- key takeaways}
\author{Gelman et al.\ 2020 \;|\; arXiv:2011.01808}
\date{}

\begin{document}
\maketitle

\section*{Inference vs.\ workflow}

\textbf{Bayesian inference} -- the math:
$p(\theta\mid y) \propto p(\theta)\,p(y\mid\theta)$. One line.

\textbf{Bayesian workflow} -- includes the three steps of
model building, inference, and model checking/improvement,
along with the comparison of different models.

\section*{Why we need a workflow}

\begin{enumerate}
    \item Computation is hard; you have to try simpler models,
          approximations, and exploratory fits.
    \item You don't know the right model up front; you expand it as
          the data and questions develop.
    \item Even with perfect data and a good model, understanding it
          requires comparing it to related models.
    \item Different models give different conclusions; reporting
          several is the only honest way to convey model uncertainty.
\end{enumerate}

\noindent Bayesian is a process, not a one-shot.

\section*{Empirical Bayes vs.\ multilevel model}

Multilevel modelling is the \emph{formalization of empirical Bayes}.
Both want to estimate group-level parameters $u_p$ (e.g.\ each
participant's deviation from the group mean), but they differ in how
they treat the prior variance $\sigma_u^2$.

\textbf{Empirical Bayes (older, two-step):}
\begin{enumerate}
    \item Fit each group separately, obtain $\hat u_p$.
    \item Estimate $\hat\sigma_u^2 = \mathrm{Var}(\hat u_p)$ from those
          estimates.
    \item Plug $\hat\sigma_u$ in as a \emph{fixed} prior and shrink the
          $\hat u_p$ toward $0$ using a James-Stein-style formula.
\end{enumerate}

\textbf{Multilevel Bayes (now, one-step):}
$\sigma_u$ itself has a hyperprior; everything is estimated jointly:
\[
    u_p \sim \mathcal{N}(0, \sigma_u^2),
    \qquad
    \sigma_u \sim \text{StudentT}(3, 0, 1).
\]

\medskip
\noindent In the BGAM model, the term
\texttt{(1 + condition | participant)} is exactly this construction:
each participant gets their own intercept and condition slope, and the
variance of those random effects is estimated jointly with everything
else --- one line replaces what used to be a multi-step empirical
Bayes procedure.

\section*{Start simple, add features gradually}

\textbf{Template vs.\ ladder --- two principles at different levels}.
Section 2.1 of Gelman et al.\ recommends starting from a published
\emph{template} (their model serves as the conceptual target).
Section 5.2 recommends fitting via a \emph{ladder} of simpler
intermediate models. These are complementary:

\begin{itemize}\setlength{\itemsep}{0pt}
    \item \textbf{Template} = the final target spec (e.g.\ Nalborczyk
          BGAM + my 3D extension). Tells me \emph{what} to build.
    \item \textbf{Ladder} = incremental fits walking up to that spec.
          Tells me \emph{how} to build it without losing track of
          where problems originate.
\end{itemize}

\noindent My earlier mistake was jumping straight from the template
into v4\_k12 (all features at once); without the ladder I cannot
attribute the $\sim$37\% amplitude shrinkage to any single feature.

\medskip
\noindent Gelman's general advice for complex models: start with a
model missing some important features and gradually add them in.

BGAM hits all three of Gelman's ``hard problem'' markers:
correlations (in random effects and in the t2 basis), hierarchical
structure (participants), and time-varying parameters (the time
dimension of t2). Plus a measurement-error model on top.

Going straight to the full v5\_k20 means that if anything is wrong
(e.g.\ the $\sim 37\%$ amplitude shrinkage), there are too many
suspects.

\textbf{Remedy: fit a ladder of models}, each adding one feature,
using VI (\texttt{algorithm = "meanfield"}) for fast diagnostics:

\begin{enumerate}\setlength{\itemsep}{0pt}
    \item \texttt{y \textasciitilde\ condition}
    \item \texttt{+ (1 | participant)}
    \item \texttt{+ (1 + condition | participant)}
    \item \texttt{+ s(time, by = condition)}
    \item \texttt{+ s(xproj, yproj, by = condition)}
    \item \texttt{+ t2(xproj, yproj, time, by = condition)}
    \item \texttt{+ | se(eeg\_sd)}
\end{enumerate}

\noindent Compare amplitudes / diagnostics step by step; the feature
that first causes shrinkage or divergence is the one to investigate.
Run HMC only on the final spec and a couple of intermediate ones for
verification.

\section*{Shortcuts are legitimate}

Real workflow is disorderly: attention bounces between data exploration,
theory, computing, and interpretation. Time and compute force every
practitioner to take shortcuts.

Key Gelman point: \emph{not} fitting a model at all is usually worse
than fitting it with an approximate computation. Shortcuts are fine
if you treat them as deliberate \textbf{approximations to the full
workflow} and document them so readers can interpret the results in
context.

\medskip
\noindent My BGAM project will rely on several documented shortcuts:
\begin{itemize}\setlength{\itemsep}{0pt}
    \item HMC compute capped at $\sim$24 h on a 16-core HPC node.
    \item Monte Carlo size 30 instead of the 500 used by Nalborczyk
          et al.\ (2025).
    \item Prior predictive checking skipped (using brms defaults).
    \item Real M/EEG application left as future / PhD work.
\end{itemize}

\noindent Listing these in the methods section is honest research,
not weakness.

\section*{Default priors are data-informed (not ``no prior'')}

When I removed the custom priors from BGAM (on Nalborczyk's advice),
brms didn't drop priors entirely --- it switched to its own defaults,
most of which are \textbf{weakly informative} with scale derived from
the response variable:

\begin{itemize}\setlength{\itemsep}{0pt}
    \item \texttt{b} (regression coefficients): improper flat prior
          (truly uninformative).
    \item \texttt{sd} (random-effect SDs) and \texttt{sds} (smooth SDs):
          \texttt{student\_t(3, 0, max(2.5, SD of response))}.
    \item \texttt{Intercept}: \texttt{student\_t(3, median(y), max(2.5, SD of y))}.
\end{itemize}

This is essentially empirical Bayes: the prior scale is estimated
from the data instead of being fixed by hand. It can be viewed as an
approximation to a fully hierarchical model in which that scale
parameter would itself have a hyperprior.

\medskip
\noindent \textbf{Why brms does this}:
\begin{enumerate}\setlength{\itemsep}{0pt}
    \item \emph{Scale invariance} -- prior behaves consistently no
          matter the unit of the response ($\mu$V vs.\ mV vs.\ V).
    \item \emph{Generic applicability} -- one default works across
          EEG, income, reaction times, etc.
    \item \emph{Avoids catastrophic priors} -- flat priors can cause
          improper posteriors and MCMC instability.
    \item \emph{Computational shortcut} -- skips one level of
          hierarchy that would otherwise require sampling.
\end{enumerate}

\medskip
\noindent \textbf{Thesis note} (to write later):

\begin{quote}\small
We adopted weakly informative default priors as implemented in brms
(B\"urkner 2017), following Gelman et al.\ (2008). For random-effect
and smooth-term SDs, this corresponds to a Student-t($3, 0,$ scale)
with the scale derived from the response variable -- a computationally
efficient empirical Bayes approximation to a fully hierarchical prior
specification.
\end{quote}

\section*{Prior predictive check on the BGAM (run locally, 2026-05-12)}

Ran \texttt{brm(\ldots, sample\_prior = "only")} with the v5\_k20 spec
plus a proper \texttt{normal(0, 5)} on \texttt{b} (the flat default
cannot be sampled from). Result on the marginal \texttt{eeg\_mean}:

\begin{center}
\begin{tabular}{lrrl}
\toprule
              & Real $y$       & Prior $y_{\text{rep}}$ & Ratio \\
\midrule
Range (max)   & $\sim 2~\mu$V & $\sim 45~\mu$V        & $\sim 22\times$ \\
Range (min)   & $\sim -2~\mu$V & $\sim -45~\mu$V       & $\sim 22\times$ \\
SD            & $\sim 0.5~\mu$V & $\sim 10$--$15~\mu$V & $\sim 20$--$30\times$ \\
\bottomrule
\end{tabular}
\end{center}

\noindent The brms defaults are weakly informative \emph{in a generic
sense} but not appropriate for the EEG amplitude scale --- they imply
data that is $\sim 25\times$ wider than reality.

\begin{figure}[ht]
\centering
\begin{subfigure}{0.49\textwidth}
    \centering
    \includegraphics[width=\linewidth]{../../results/figures/ppc_prior_overlay.png}
    \caption{Density: real $y$ (dark) vs.\ prior $y_{\text{rep}}$ (light).}
\end{subfigure}
\hfill
\begin{subfigure}{0.49\textwidth}
    \centering
    \includegraphics[width=\linewidth]{../../results/figures/ppc_prior_sd.png}
    \caption{Per-draw SD: real $\sim 0.5\,\mu$V vs.\ prior $\sim$10--15\,$\mu$V.}
\end{subfigure}

\vspace{0.5em}
\begin{subfigure}{0.49\textwidth}
    \centering
    \includegraphics[width=\linewidth]{../../results/figures/ppc_prior_max.png}
    \caption{Per-draw max: real $\sim 2\,\mu$V vs.\ prior $\sim$45\,$\mu$V.}
\end{subfigure}
\hfill
\begin{subfigure}{0.49\textwidth}
    \centering
    \includegraphics[width=\linewidth]{../../results/figures/ppc_prior_min.png}
    \caption{Per-draw min: real $\sim -2\,\mu$V vs.\ prior $\sim -45\,\mu$V.}
\end{subfigure}
\caption{Prior predictive check for the v5\_k20 spec with brms defaults
plus \texttt{normal(0, 5)} on \texttt{b}. The dark line/bar is the real
data statistic; the light histogram is the distribution under the prior
across simulated datasets.}
\end{figure}

\section*{Why the supervisor said ``drop the priors'' (and why it's not brute-force)}

Reading the supervisor's recommendation as ``trust the data, distrust
the prior'' is partly right but understates the design. The actual
logic is:

\begin{enumerate}\setlength{\itemsep}{0pt}
    \item With $n = 31{,}110$ observations and a measurement-error
          term, the likelihood carries strong information.
    \item Under weak priors, the posterior is dominated by the
          likelihood; the prior contributes little.
    \item Therefore, if amplitude shrinkage persists under weak
          priors, the prior is \emph{ruled out} as the source ---
          the shrinkage must come from the smoothing penalty, the
          random-effects shrinkage, or the measurement-error model.
\end{enumerate}

\noindent So switching to default priors is best read as a
\textbf{controlled experiment / diagnostic}, not as ``brute force.''
The PPC confirms that the weak prior really is dominated by the data
(implied range $\sim 25\times$ too wide $\Rightarrow$ no
prior-induced shrinkage toward 0).

\medskip
\noindent Important caveat: this strategy only works because (a) the
data are informative, (b) the likelihood geometry is well-behaved,
and (c) PPC has been performed to check that the weak prior does not
imply absurd predictions. With $n = 5$ participants and a small
dataset, weaker priors could otherwise allow the posterior to wander.

\section*{\S 3.2 -- 3.3: fast-and-approximate strategies}

These two sections together formalise the ``fit fast, fail fast''
slogan. They tell you (a) how long to run the sampler and (b)
which approximation to use, depending on what each fit is for.

\subsection*{\S 3.2 Run length per goal}

Sampling effort should match the purpose of the fit:

\begin{itemize}\setlength{\itemsep}{0pt}
    \item \textbf{Debug}: 100s of iter -- just enough to see the
          model compile and not blow up.
    \item \textbf{Ladder diagnostic}: VI / Pathfinder / short HMC --
          you need the direction, not the digits.
    \item \textbf{Final single fit (paper)}: 2--3k iter, R-hat
          $< 1.01$, ESS $> 400$, 0 divergences.
    \item \textbf{Many-fit Monte Carlo evaluation}: minimum-but-useful
          per fit; MC averaging over sims smooths individual noise.
    \item \textbf{Tail / rare-event inference}: very long; ESS for
          quantiles needs more samples than for the posterior mean.
\end{itemize}

\noindent In my projects:

\begin{itemize}\setlength{\itemsep}{0pt}
    \item BGAM v5\_k20 HMC: 4 chains $\times$ 3000 iter, 27.5 h --
          final fit, conservative.
    \item BGAM ladder (A/B/C/D): VI 30 min each, used only for
          \emph{direction-finding}; Pareto k flagged as caveat.
    \item Stage 2 OU coverage MC: 2 chains $\times$ 1500 iter $\times$
          30 sims -- light per fit, MC averaging carries the
          aggregate accuracy.
\end{itemize}

\subsection*{\S 3.3 Approximate algorithm vs.\ approximate model}

Gelman distinguishes two failure modes of ``cheap'' inference:

\begin{itemize}\setlength{\itemsep}{0pt}
    \item \textbf{Algorithm approximation} -- same model, simpler
          sampler. VI (mean-field, full-rank), Pathfinder, INLA,
          Laplace, EP, Whittle likelihood. Fast, but can fail
          silently (e.g.\ posterior SD underestimation).
    \item \textbf{Model approximation} -- simpler model fit
          exactly. Dropping random effects, dropping ME term,
          replacing a hierarchical model by a pooled one. Loses
          modelling capacity in exchange for tractability.
\end{itemize}

\noindent The key warning is that approximation \emph{fails
silently}: HMC reports divergences and high R-hat, but VI may
return a number that looks fine while being wrong. Diagnostics:

\begin{itemize}\setlength{\itemsep}{0pt}
    \item VI: Pareto $k$ (should be $< 0.7$; mine was 11.5 for the
          A-only run, 3.1 for the merged run -- treat the SD/CI
          columns as approximate).
    \item INLA / Laplace: check marginal vs.\ joint posterior shape.
    \item Whittle: compare against time-domain on a short subset
          before committing.
\end{itemize}

\subsection*{Where each appeared in my workflow so far}

\begin{center}
\begin{tabular}{lll}
\toprule
Project & Approximation used & Type \\
\midrule
BGAM ladder (A/B/C/D) & VI (mean-field), 20k iter        & algorithm \\
BGAM ablation A       & drop \texttt{| se(eeg\_sd)}      & model \\
BGAM Monte Carlo      & 30 sims (not 500)                & MC budget \\
Stage 2 OU (single)   & none                              & --- \\
Stage 2 OU coverage   & 2 chains, 1500 iter, 30 sims      & run length \\
Stage 3 (planned)     & Whittle likelihood at $T\ge 10^3$ & algorithm \\
\bottomrule
\end{tabular}
\end{center}

\noindent The unifying rule: \emph{use the cheap option to find the
direction; verify with the expensive option before reporting}. The
Bayesian workflow is not ``pick one sampler and run forever'', it
is a deliberate sequence of approximations with each one's failure
mode documented.

\section*{Stage 2 OU: $\sigma^2$ prior is slightly too wide}

The Monte Carlo evaluation of the OU model (30 sims $\times$
$T \in \{50, 100, 200\}$, truth $(\alpha,\sigma^2) = (1, 1)$)
showed a consistent positive bias for $\sigma^2$:

\begin{center}
\begin{tabular}{lrrrr}
\toprule
$T$ & $\sigma^2$ bias & $\sigma^2$ var & $\sigma^2$ MSE & 95\% coverage \\
\midrule
50  & $+0.54$ & 0.12 & 0.41 & 100\% \\
100 & $+0.36$ & 0.21 & 0.34 & 96.7\% \\
200 & $+0.23$ & 0.22 & 0.27 & 83.3\% \\
\bottomrule
\end{tabular}
\end{center}

\noindent The bias decreases with $T$ (consistent estimator), but
the 95\% coverage drops below the nominal level at $T = 200$. The
prior predictive check explains why: with $\sigma^2 \sim
\mathrm{Half\text{-}t}(3, 0, 2.5)$ the implied prior median of
$\sigma^2$ is $\approx 1.92$, almost twice the truth, so when the
likelihood is weak the posterior is pulled upward.

\subsection*{v2 fix and the parameter-coupling trade-off}

I tightened the prior to $\sigma^2 \sim \mathrm{Half\text{-}t}(3, 0, 1)$
(prior median $\approx 0.77$), re-ran the MC evaluation, and got:

\begin{center}\small
\begin{tabular}{lrrrrrr}
\toprule
$T$ & $\sigma^2$ bias (v1 / v2) & $\sigma^2$ MSE (v1 / v2)
    & $\sigma^2$ cov (v1 / v2) & $\alpha$ bias (v1 / v2) \\
\midrule
50  & +0.54 / +0.12 & 0.41 / 0.08 & 100\% / 100\%   & +0.03 / +0.26 \\
100 & +0.36 / +0.10 & 0.34 / 0.12 & 96.7\% / 96.7\% & +0.09 / +0.21 \\
200 & +0.23 / +0.07 & 0.27 / 0.13 & 83.3\% / 93.3\% & +0.03 / +0.11 \\
\bottomrule
\end{tabular}
\end{center}

\noindent $\sigma^2$ recovery improves substantially (bias and MSE
roughly halved, coverage at $T = 200$ goes from $83\%$ back up to
$93\%$). But \emph{$\alpha$ bias gets worse} because the posterior
$(\alpha, \sigma^2)$ is positively correlated: pulling $\sigma^2$
down via the prior pushes $\alpha$ up to keep the likelihood near
its peak. This is the classical $\sigma^2 / \alpha$ identifiability
issue in OU.

\textbf{Lesson:} a prior tweak that fixes one parameter's bias
typically leaks into correlated parameters. Robust fixes have to
target the joint geometry, either by:

\begin{itemize}\setlength{\itemsep}{0pt}
    \item placing the prior on a less-correlated transform of the
          parameters (e.g.\ $\log\sigma^2 - \log\alpha$), or
    \item using a non-centred / Cholesky parametrisation that
          breaks the geometry, or
    \item accepting the trade-off and reporting both versions as
          a sensitivity analysis (\S 6.4).
\end{itemize}

\noindent This is itself a finding worth recording: the workflow
caught a coupling that a single-fit run would have missed.

\section*{How to pick weakly informative priors (4-step recipe)}

When moving from simulation to real EMG data, the same OU model
needs a new prior. The principle is \emph{never} centre the prior
on the unknown truth -- pick from domain knowledge:

\begin{enumerate}\setlength{\itemsep}{0pt}
    \item \textbf{Support} -- what range can the parameter take?
          ($\alpha > 0$, $\sigma^2 > 0$.)
    \item \textbf{Scale} -- what is the plausible order of
          magnitude from domain knowledge?
    \item \textbf{Family} -- positive + heavy-tailed: Lognormal,
          Half-Student-t, Half-Cauchy.
    \item \textbf{Prior predictive check} -- simulate fake data
          and verify it is physically reasonable.
\end{enumerate}

\subsection*{Concrete example: priors for real resting EMG}

\begin{itemize}\setlength{\itemsep}{0pt}
    \item \textbf{$\alpha$} -- resting EMG autocorrelation decays on
          the millisecond scale, so the decorrelation time
          $1/\alpha$ is roughly $1$--$100$\,ms, giving
          $\alpha \sim 10$--$1000$\,rad/s.\\
          Prior: $\alpha \sim \mathrm{Lognormal}(\log 100,\, 1)$
          $\Rightarrow$ median $100$, 95\% CI roughly
          $[14,\,720]$\,rad/s.
    \item \textbf{$\sigma^2$} -- resting EMG RMS is around
          $5\,\mu$V, so $\sigma^2 \sim 25\,\mu\mathrm{V}^2$.\\
          Prior: $\sigma^2 \sim \mathrm{Half\text{-}t}(3,\, 0,\, 25)$
          $\Rightarrow$ median $\sim 19$, heavy-tailed enough to
          accommodate noisy electrodes.
\end{itemize}

\noindent Neither prior is centred on a true value -- both are
chosen from the \emph{order of magnitude} suggested by the EMG
literature and pilot recordings. Prior PC on simulated paths then
confirms that the implied amplitude and timescale match what real
resting EMG looks like.

\end{document}
